% !PRINTED 162
\weberpage{162}
%<<CONT
to take an arbitrary form $\varphi(x)$ of the $n$ variables $x$,
and to form the functions:
\weq{2}{\varphi[A(x)],\ \varphi[B(x)],\ \varphi[C(x)],\ \ldots}
for all substitutions of the group, and to take any
symmetric function of the forms (2) for $\Phi$.

For if one applies to $(x)$ a substitution of the group $S$,
the totality of the functions (2) does not change; only
their order becomes a different one.

One can conceive the notion of invariants still more generally,
as follows:

2. \gesperrt{A form $\mathit\Psi(x_1,x_2,\ldots,x_n)$ is called an
invariant of the group $S$ also when it takes on
constant factors when the substitutions
$A,B,C \ldots$ of the group $S$ are applied to the
variables $(x)$.}

This property is expressed by formulae as follows:
\weq{3}{\mathit\Psi[A(x)]=\alpha\,\mathit\Psi(x),\ \mathit\Psi[B(x)]=\beta\,\mathit\Psi(x),\ \mathit\Psi[C(x)]=\gamma\,\mathit\Psi(x)\ldots}
in which the coefficients $\alpha,\beta,\gamma\ldots$ are independent of the $x$.

When a distinction is necessary, we shall call the forms
defined in 1. \gesperrt{absolute invariants}, and those defined in 2. \gesperrt{relative
invariants} of the group $S$.

From the formulae (3) there results:
\weq{4}{\mathit\Psi[AB(x)] = \alpha\beta\,\mathit\Psi(x),}
and from this it follows that the factors $\alpha,\beta,\gamma,\ldots$ in their totality,
under composition by multiplication, must form a
group.

Between this (commutative) group and the group $S$
there is in general a multiple isomorphism, since
different substitutions of $S$ can lead to the same factor $\alpha$.

If $\mu$ is the degree of the group $S$, then the degree of each
element $A, B \ldots$ of $S$ is a divisor of $\mu$, and consequently
$A^\mu = B^\mu \cdots$ is equal to the identical substitution.

From this it follows that the factors $\alpha,\beta,\gamma\ldots$ are \gesperrt{$\mu$th roots of unity}.

If $e$ is the least positive number satisfying the condition
\weq{5}{\alpha^e = \beta^e = \gamma^e = \cdots = 1}
then the $\alpha, \beta, \gamma, \ldots$ are at the same time $e$th roots
%CONT>>
